\documentclass[aps,prd,twocolumn,superscriptaddress,nofootinbib,showkeys]{revtex4-2}
\usepackage{amsmath,amssymb,hyperref}
\begin{document}
\title{CERN+1: An Independent LHC Toolkit That Checks Itself --- Wino Lifetimes, Two Reproduced ATLAS Limits, and a Model-versus-Collider Panel}
\author{Danny Eldridge}
\thanks{CERN+1 is independent work. It is not affiliated with, sponsored by or endorsed by CERN or any LHC experiment; the name means ``the chain, plus one more check''.}
\email{dannyeldridge54@gmail.com}
\affiliation{Independent researcher, Oklahoma City, Oklahoma, USA}
\date{18 November 2026}
\begin{abstract}
CERN+1 is a small, open toolkit that repeats the main steps an LHC search uses to test a wino-like dark-matter candidate, with an independent check at each step. Two cross-section routes (a spline through the Resummino tables and Drell--Yan scaling, $\tau_0=4m^2/s$) agree to 1\% at 14~TeV; this is interpolation and route spread only, not the PDF and scale uncertainty of the tables. A mass-dependent pure-wino lifetime engine gives $\delta m=164.1$~MeV and $c\tau=5.83$~cm at 1~TeV. A benchmark-anchored disappearing-track recast, with a chargino boost distribution computed from the same tables, gives 661.5~GeV for ATLAS 136~fb$^{-1}$ (published 660) and 446.6~GeV for ATLAS 36.1~fb$^{-1}$ (published 460); because each yield is anchored on ATLAS's own benchmark count, this tests how the yield falls with mass. Exact Poisson CLs and pyhf differ by up to 8\% at low counts. Our independent recast reproduces the lower edge of the post-hoc SERAPIS wino window (0.67--1.85~TeV), which was itself set from an ATLAS limit. The HL-LHC reaches 975~GeV without a background systematic and 896~GeV with a 20\% one, covering 26\% or 19\% of the window. Future-collider reaches are 0.97/0.80~TeV (HL-LHC; published 0.9/0.5), 1.68/1.36~TeV (HE-LHC; 2.1/1.5) and 3.74/2.88~TeV (FCC-hh; 6.5/4.5). Two cross-section self-tests pass at 0.1\%; a charge-asymmetry test and the high-background FCC-hh comparison failed as locked and are reported. Every step was pre-registered.
\end{abstract}
\keywords{LHC; wino; chargino lifetime; disappearing track; recast; reinterpretation; CLs; HL-LHC; dark matter; reproducibility; pre-registration}
\maketitle

\section{Why a second, self-checking chain}
The cross-section module was released under the working name Janus-LHC~\cite{P1}; the toolkit is now called CERN+1.
A wino-like electroweakino is one of the cleanest dark-matter targets at a hadron collider. Its charged partner is about 160~MeV heavier than the neutral one. It lives long enough to leave a few centimetres of track that then ``disappear''. To test a specific model one needs four things: the production rate at the collider energy, the charged-state lifetime, the search's efficiency and limit, and a statistics engine. The experiments' chain is MadGraph/Pythia/Delphes for simulation, Resummino/Prospino for cross sections, and SModelS, CheckMATE or full likelihoods for reinterpretation. That chain is accurate, but it is slow, and it often fixes key inputs to one benchmark. CERN+1 recomputes each input independently and attaches a measured error to every answer.

\section{The process}
Five rules were applied to every component.
\begin{enumerate}
\item \emph{Lock before computing.} The expected outcome and the pass criterion are written, hashed and published first~\cite{P4}.
\item \emph{Two routes.} A number is trusted only when two independent methods agree, and their spread is the quoted error.
\item \emph{Reproduce before predicting.} A recast is used for predictions only after it reproduces the experiment's own limit.
\item \emph{Replace benchmarks with physics.} Fixed benchmark inputs are replaced by quantities computed from the model.
\item \emph{Report failures.} One sub-prediction below failed, and it is recorded as a failure.
\end{enumerate}

\section{Cross sections (summary of Ref.~\cite{P1})}
Route~1 is a spline of $\ln\sigma$ against $\ln m$ through the Resummino tables~\cite{Res}. Left-out points are reproduced with a median error of 0.016\% and a maximum of 0.30\%. Route~2 is Drell--Yan scaling, $s\sigma=\Phi(\tau_0)(\sqrt s/13\,\mathrm{TeV})^\kappa$ with $\tau_0=4m^2/s$ and $\Phi$ a universal function read off the 13~TeV table alone. It predicts every 13.6~TeV entry to $\le0.29\%$. The exponent is $\kappa=-0.163\pm0.002$ across the three wino channels. At 14~TeV the two routes agree to $\le1\%$. Both routes share the same 13 and 13.6~TeV inputs, so this spread does not include the PDF and scale uncertainty of the tables, which exceeds 10\% above about 1.5~TeV.

\section{Lifetime engine}
We use the two-loop fit to the wino mass splitting~\cite{IMS} and the widths $\Gamma(\tilde\chi^\pm\to\tilde\chi^0\pi^\pm)$ and $\Gamma(e\nu)$ given there. We add the $\mu\nu$ mode with the standard muon-decay phase-space function in the heavy-parent limit. Results:
\begin{itemize}
\item $\delta m$ rises from 149.5~MeV at 100~GeV to 164.5~MeV at 3~TeV.
\item $c\tau$ falls from 11.3 to 5.75~cm.
\item BR$(\pi)\approx97.9\%$ above 300~GeV (97.5\% at 100~GeV).
\item The lifetime is 0.197~ns at 666~GeV, compared with the 0.2~ns benchmark used by ATLAS.
\end{itemize}
The pre-registered window for $\delta m(1\,\mathrm{TeV})$, 160--163~MeV, was 1~MeV too low. The value is 164.1~MeV, so that sub-prediction failed. The other predictions passed.

\section{Disappearing-track recast, validated twice}
The signal yield is anchored on each search's published benchmark and scaled with our cross sections and our lifetime:
\begin{equation}
N(m)=N_{\rm b}\,\frac{\sigma(m)}{\sigma(m_{\rm b})}\,\frac{P(m)}{P(m_{\rm b})}.
\end{equation}
Here $N_{\rm b}$ is ATLAS's own expected signal count at the benchmark mass $m_{\rm b}$, and $P$ is the probability to decay between the fourth pixel layer (122.5~mm) and the first SCT layer (299~mm), at transverse boost $\beta\gamma=p_T/m$. Efficiency and acceptance are held constant in mass; only $\sigma$ and $P$ vary. $N(m)$ is compared with ATLAS's published model-independent limit $N_{95}$ in the electroweak signal region (we use their values, not ours). The test therefore checks the mass fall-off of the yield between the benchmark and the limit, not the absolute efficiency.
\begin{itemize}
\item ATLAS 136~fb$^{-1}$~\cite{AT136}: benchmark 8.6 events at 600~GeV, $N_{95}=5.03$. The limit is 666.6~GeV against the published 660~GeV ($+1.0\%$).
\item ATLAS 36.1~fb$^{-1}$~\cite{AT36}: benchmark 13.5 events at 400~GeV, $N_{95}=7.94$. The limit is 447.8~GeV against 460~GeV ($-2.7\%$).
\end{itemize}
Nothing was retuned between the two searches. With a fixed $\beta\gamma=0.5$ the limits are 666.6 and 447.8~GeV. Replacing it with a boost distribution (a parton luminosity fitted to the table's own mass dependence to 3.2\%, $1+\cos^2\theta$ production; median $\beta\gamma=0.47$--0.54) gives 661.5~GeV ($+0.2\%$) and 446.6~GeV ($-2.9\%$); this refinement was locked first (D1307). A like-for-like HL-LHC extrapolation (14~TeV, 3000~fb$^{-1}$, the Run-2 background $3.0\pm0.7$ scaled with luminosity to $b=66.2$, Run-2 efficiency, ITk not modelled) reaches 975~GeV with exact-Poisson $N_{95}=17.6$ and no background systematic, and 896~GeV with a 20\% background uncertainty (pyhf, $N_{95}=28.6$).

\section{Statistics}
pyhf and exact Poisson CLs disagree by up to 8\% at low counts. For example, at $b=n=3$ they give 5.00 and 5.40 signal events. The toolkit reports the exact value whenever counts are small. The exact route has no background uncertainty; pyhf is used when one is needed.

\section{Model versus collider}
The SERAPIS model has a post-hoc thermal-wino window of 0.67--1.85~TeV, making up 5--41\% of the dark matter~\cite{SE}. Its lower edge was set from the ATLAS limit~\cite{AT26}, so it is not a prediction; our independent recast reproduces that edge. The panel reads the window from the model's theory map and overlays the CERN+1 reach. The HL-LHC covers 670--975~GeV (26\% of the window) without a background systematic, or 670--896~GeV (19\%) with a 20\% one. The rest needs future colliders (below) or a gamma-ray line search.

\section{Symmetry self-tests}
The proton contains more $u$ than $d$ quarks, and Drell--Yan production depends on energy mainly through $\tau_0$. Three self-tests therefore run on every cross-section table.
\begin{itemize}
\item (T1) The charge asymmetry $A=\sigma(\tilde\chi_2^0\tilde\chi_1^+)/\sigma(\tilde\chi_2^0\tilde\chi_1^-)$ should exceed 1 and rise with mass.
\item (T2) At equal $\tau_0$, $A$ should be the same at 13 and 13.6~TeV.
\item (T3) The ratio $R=[\sigma(\tilde\chi_2^0\tilde\chi_1^+)+\sigma(\tilde\chi_2^0\tilde\chi_1^-)]/\sigma(\tilde\chi_1^+\tilde\chi_1^-)$ should depend on $\tau_0$ only.
\end{itemize}
T2 and T3 pass to 0.093\% and 0.050\%. T1 failed as locked. $A$ rises from 1.58 at 100~GeV to 3.36 at 1.5~TeV, then falls to 3.28 at 2~TeV. The turnover sits where the table's PDF uncertainty becomes very large. The self-test now marks that region as PDF-limited, and T2 and T3 guard the tables against corruption.

\section{Future colliders}
The module reuses the Run-2 recast efficiency and decay window. It takes cross sections from the Drell--Yan scaling route, an extrapolation far beyond 13.6~TeV that we label rough. It scales the background with luminosity, times a factor of 1 (low) or 10 (high). The published reference is Ref.~\cite{FCC}, whose two variants are $0.2\times$ and $5\times$ a central background; ours are $1\times$ and $10\times$ a luminosity-scaled Run-2 background, so the bands are not identical. Results, ours against published:
\begin{itemize}
\item HL-LHC (14~TeV, 3~ab$^{-1}$): 0.97 and 0.80~TeV, against 0.9 and 0.5~TeV.
\item HE-LHC (27~TeV, 15~ab$^{-1}$): 1.68 and 1.36~TeV, against 2.1 and 1.5~TeV.
\item FCC-hh (100~TeV, 30~ab$^{-1}$): 3.74 and 2.88~TeV, against 6.5 and 4.5~TeV.
\end{itemize}
The locked test compared each variant with the edge of the published band (4.5~TeV): the low-background value 3.74~TeV is 17\% short (pass) and the high-background 2.88~TeV is 36\% short (fail, criterion 25\%). Against the published low-background value (6.5~TeV) ours is 42\% short. The weak link is the cross section extrapolated to 100~TeV: with exact scaling ($\kappa=0$) the FCC-hh reach rises to 4.03 and 3.11~TeV. Every FCC-hh variant still covers the whole SERAPIS window.

\section{What the toolkit adds}
\begin{itemize}
\item Every answer carries a measured interpolation or route-spread error (the tables' own PDF and scale uncertainty is not yet propagated).
\item Lifetimes are computed as a function of mass rather than fixed at one benchmark (a 1.4\% effect at the limit).
\item Exact CLs at low counts, where the asymptotic formulae are off by up to 8\%.
\item The recast is validated on two searches before it is used.
\item Answers take microseconds to seconds, against minutes to hours for the reference codes.
\end{itemize}
There is no head-to-head comparison with SModelS or CheckMATE yet, no full event generation or detector simulation. CERN+1 uses the experiments' published efficiencies, and its cross sections rest on the reference tables until our own PDF-based leading-order route is added.

\emph{How we came up with it.}
CERN+1 grew out of the SERAPIS research, which predicts wino dark matter in a window that LHC searches can test. To score that prediction honestly we needed the experiments' numbers, fast and with errors. When a parton-distribution download was blocked, we asked what physics alone fixes. Drell--Yan scaling turned one energy's table into a predictor for every energy, and the leftover correction came out the same in every channel. The habits that built our inflation calculator --- lock first, two routes, reproduce before predicting --- were then applied to each LHC step.

\emph{Method note.} Prepared with an AI assistant (Claude, Anthropic). Code: \texttt{tools/lhc}; notes D1299--D1307 in the archive~\cite{AEGIS}.

\begin{thebibliography}{9}
\bibitem{P1} D.~Eldridge, ``Janus-LHC: Wino Pair Cross Sections at Any LHC Energy from Drell--Yan Scaling,'' Zenodo (2026), doi:10.5281/zenodo.24342227.
\bibitem{Res} J.~Fiaschi, B.~Fuks, M.~Klasen and A.~Neuwirth, ``Electroweak superpartner production at 13.6 TeV with Resummino,'' Eur.\ Phys.\ J.\ C \textbf{83}, 707 (2023), arXiv:2304.11915.
\bibitem{IMS} M.~Ibe, S.~Matsumoto and R.~Sato, ``Mass splitting between charged and neutral winos at two-loop level,'' Phys.\ Lett.\ B \textbf{721}, 252 (2013), arXiv:1212.5989.
\bibitem{AT136} ATLAS Collaboration, ``Search for long-lived charginos based on a disappearing-track signature using 136~fb$^{-1}$,'' Eur.\ Phys.\ J.\ C \textbf{82}, 606 (2022), arXiv:2201.02472.
\bibitem{AT36} ATLAS Collaboration, ``Search for long-lived charginos based on a disappearing-track signature in pp collisions at $\sqrt s=13$~TeV with the ATLAS detector,'' JHEP \textbf{06}, 022 (2018), arXiv:1712.02118.
\bibitem{SE} D.~Eldridge, ``SERAPIS-$E$ v2,'' Zenodo (2026), doi:10.5281/zenodo.24340373.
\bibitem{P4} D.~Eldridge, ``Lock Before You Compute,'' Zenodo (2026), doi:10.5281/zenodo.24291026.
\bibitem{FCC} T.~Han, S.~Mukhopadhyay and X.~Wang, ``Electroweak Dark Matter at Future Hadron Colliders,'' Phys.\ Rev.\ D \textbf{98}, 035026 (2018), arXiv:1805.00015.
\bibitem{AT26} ATLAS Collaboration, arXiv:2603.08315.
\bibitem{AEGIS} Archive: github.com/dannyeldridge54/aegis.
\end{thebibliography}
\end{document}
